\chapter*{คำนำ}
\addcontentsline{toc}{chapter}{คำนำ}

ความก้าวหน้าอย่างก้าวกระโดดของเทคโนโลยีความเป็นจริงขยาย (Extended Reality หรือ XR) ซึ่งครอบคลุมความเป็นจริงเสมือน (Virtual Reality: VR) ความเป็นจริงเสริม (Augmented Reality: AR) และความเป็นจริงผสม (Mixed Reality: MR) ได้เข้ามามีบทบาทอย่างยิ่งในการปฏิรูปกระบวนการจัดการเรียนรู้วิทยาศาสตร์ เทคโนโลยี วิศวกรรมศาสตร์ และคณิตศาสตร์ หรือสะเต็มศึกษา (STEM Education) ในศตวรรษที่ 21 การเรียนรู้ปรากฏการณ์ทางกายภาพที่มีความซับซ้อน เป็นนามธรรม หรือมีข้อจำกัดด้านความปลอดภัยและงบประมาณในการทดลองจริง สามารถถูกจำลองขึ้นในรูปแบบสามมิติเชิงโต้ตอบที่มีความเที่ยงตรงสูง ช่วยให้ผู้เรียนเกิดความเข้าใจเชิงมโนทัศน์อย่างลึกซึ้งผ่านประสบการณ์ตรง

ในอดีต การพัฒนาแอปพลิเคชัน XR มักเผชิญกับอุปสรรคสำคัญจากปัญหาความแตกแยกของระบบ (Ecosystem Fragmentation) เนื่องจากผู้พัฒนาต้องเขียนโค้ดแยกตามชุดคำสั่งและอุปกรณ์เฉพาะของผู้ผลิตแต่ละราย (Proprietary SDKs) อย่างไรก็ดี การถือกำเนิดขึ้นของมาตรฐานเปิด \textbf{OpenXR} โดยความร่วมมือของกลุ่มพันธมิตรระดับโลก Khronos Group ได้สร้างจุดเปลี่ยนครั้งประวัติศาสตร์ ด้วยการกำหนดมาตรฐานกลางระดับสากลที่เชื่อมประสานระหว่างซอฟต์แวร์ประยุกต์และฮาร์ดแวร์ทุกแพลตฟอร์มอย่างเป็นเอกภาพ ส่งผลให้นักการศึกษาและนักพัฒนาสามารถสร้างสรรค์ระบบจำลองการทดลองเสมือนจริงที่ทำงานข้ามแพลตฟอร์มได้อย่างไร้รอยต่อ

ตำราวิชาการเรื่อง \textbf{``การศึกษาและการจำลองเชิงสะเต็มด้วยมาตรฐาน OpenXR และการประมวลผลเชิงพื้นที่ (OpenXR and Spatial Computing for STEM Education)''} เล่มนี้ ได้รับการรังสรรค์ขึ้นเพื่อเป็นเอกสารทางวิชาการและคู่มือการเรียนรู้เชิงลึกที่เชื่อมโยงทฤษฎีทางฟิสิกส์ สถาปัตยกรรมคอมพิวเตอร์กราฟิกส์ และการคำนวณทางวิทยาศาสตร์เข้าด้วยกันอย่างเป็นระบบ เนื้อหาภายในเล่มแบ่งออกเป็น 7 บทเรียนหลัก ครอบคลุมตั้งแต่สถาปัตยกรรมพื้นฐานของ OpenXR, ระบบพิกัดและการติดตามการเคลื่อนที่ 6 องศาอิสระ (6DoF), การปฏิสัมพันธ์เชิงพื้นที่ผ่านการตรวจจับโครงร่างมือ 21 จุดร่วม (21-Joint Skeletal Hand Tracking), เอนจินฟิสิกส์และการจำลองปรากฏการณ์ทางสะเต็ม, การประมวลผลทางวิทยาศาสตร์ด้วยภาษา Python และ PyOpenXR, การบูรณาการปัญญาประดิษฐ์และคอมพิวเตอร์วิทัศน์, ตลอดจนการพัฒนาและประยุกต์ใช้ห้องปฏิบัติการเสมือนจริงขั้นสูง

ผู้เขียนหวังเป็นอย่างยิ่งว่า ตำราเล่มนี้จะเป็นประโยชน์ต่อนักศึกษา คณาจารย์ นักวิจัย ตลอดจนผู้สนใจในการประยุกต์ใช้เทคโนโลยีความเป็นจริงเสมือนเพื่อการศึกษา และเป็นส่วนหนึ่งในการขับเคลื่อนการจัดการศึกษาสะเต็มเชิงรุกให้ก้าวสู่ระดับสากลอย่างยั่งยืน

\vspace{1.5cm}
\begin{flushright}
    \textbf{ผู้ช่วยศาสตราจารย์ ดร.ชีวะ ทัศนา}\\[4pt]
    สาขาวิชาฟิสิกส์ คณะวิทยาศาสตร์และเทคโนโลยี\\[2pt]
    มหาวิทยาลัยราชภัฏรำไพพรรณี
\end{flushright}
\clearpage
